% R008 — Guide to Cultivating Complex Analysis (Bahasa Indonesia)
% Reader-facing translation unit: source/ca-v1.9/ca.tex lines 556–667.
% Source release: v1.9 (commit a4d25d9ba37a486a5a020219eb8bd4e6602fd14c).
% Translation provenance: OpenAI Codex gpt-5.6-sol, Ultra, acting on the user's request.
% Preserve the source labels, cross-reference macros, bibliography keys, and
% diagram structure when this unit is integrated into the full driver.

\chapter*{Pendahuluan} \label{ch:intro}
\addcontentsline{toc}{chapter}{Pendahuluan}
\markboth{PENDAHULUAN}{PENDAHULUAN}

\begin{myepigraph}
Jika Anda tidak dapat membuktikan bahwa seseorang keliru, jangan panik. Anda selalu dapat memanggilnya dengan nama-nama buruk.

---Oscar Wilde
\end{myepigraph}

Tujuan buku ini adalah mengajarkan mata kuliah analisis kompleks tingkat
pascasarjana selama satu semester bagi mahasiswa pascasarjana yang baru masuk.\footnote{Saya menulisnya khusus untuk mengajar Math 5283 di Oklahoma State University.}
Tujuh bab pertama merupakan rangkaian semester pertama yang wajar dalam
urutan dua semester; semester kedua dapat membahas beberapa variabel kompleks
(misalnya, \cite{scv:book}) atau mungkin analisis harmonik.
Buku ini mungkin juga digunakan untuk rangkaian dua semester yang lebih
elementer jika lampiran dibahas terlebih dahulu dan semua bagian opsional dari
teks utama juga dibahas.
Kami mengasumsikan pengetahuan dasar analisis tingkat sarjana pada variabel
real, yang di beberapa perguruan tinggi disebut \myquote{kalkulus lanjut}.
Teks ini mengasumsikan pengetahuan tentang ruang metrik dan kalkulus
diferensial beberapa variabel, tetapi jika pembaca belum menguasai topik-topik
tersebut atau belum pernah mempelajarinya, hasil-hasil yang diperlukan
disajikan (beserta pembuktiannya) dalam lampiran.
Dengan demikian, prasyarat dasar untuk mata kuliah ini setidaknya satu
semester analisis tingkat sarjana jika lampiran juga dibahas atau dibaca; jika
mahasiswa telah mempelajari ruang metrik dan pemetaan di $\R^2$, kuliah dapat
langsung dimulai pada \Chapterref{ch:complexplane}.
Pengetahuan yang sangat dasar tentang aljabar linear dan aljabar abstrak
tingkat sarjana juga bermanfaat.

Prasyarat analisis terutama dapat ditemukan dalam
\cites{ra:book,ra:book2,Rudin:principles}.
Bacaan lanjutan yang direkomendasikan tentang analisis kompleks adalah
\cites{Boas,Conway1,Conway2,Rudin,Ullrich}.
Lihat bab yang dengan tepat diberi nama
\hyperref[ch:furtherreading]{Bacaan Lanjutan}.

Buku ini berpandangan bahwa kita tidak perlu mendefinisikan ulang dan
membuktikan kembali hal-hal yang telah dibahas dalam mata kuliah analisis real
tingkat sarjana, terutama yang berkaitan dengan pemetaan bidang.
Kita dapat segera beralih ke fungsi holomorf sebagai solusi persamaan
Cauchy--Riemann, misalnya.
Hubungannya adalah memahami baik turunan suatu pemetaan bidang maupun
perkalian dengan bilangan kompleks sebagai matriks real berukuran $2 \times 2$.
Ketika memperkenalkan integral garis, kami menghubungkannya dengan integral
garis yang telah dipelajari mahasiswa dalam kalkulus.
Teorema fungsi invers holomorf dapat diperkenalkan sejak awal sebagai
konsekuensi teorema fungsi invers standar di $\R^2$.
Garis besar pembuktian analisis kompleks murni ditunda ke latihan.
Ini bukan sekadar langkah penghematan waktu.
Intinya adalah menekankan bahwa kita tidak sedang mendefinisikan dunia yang
sama sekali baru dan berbeda.

Kami juga berusaha memperkenalkan pendekatan $z,\bar{z}$ alih-alih hanya
pendekatan $x,y$ murni.
Misalnya, kami memperkenalkan dan menggunakan operator Wirtinger.
Ini benar-benar cara yang lebih baik untuk memikirkan variabel kompleks.

Kami berusaha untuk tidak mendefinisikan terminologi atau notasi yang
bertentangan dengan apa yang telah dipelajari pembaca sebelumnya.
Terutama, istilah \myquote{diferensiabel} umumnya kami sisakan untuk konteks
turunan real, dan kami menggunakan \myquote{diferensiabel kompleks} bila
diperlukan.
Meskipun demikian, kami sering menulis \myquote{(real) diferensiabel} atau
\myquote{diferensiabel (dalam arti real)} untuk memperjelas bahwa yang
dimaksud adalah diferensiabilitas real.

Terakhir, beberapa bagian dalam tujuh bab pertama diberi tanda $\star$ dan
dapat dilewati dengan mudah pada pembacaan pertama (meskipun ini bukan berarti
bagian tersebut tidak penting, hanya saja tidak diperlukan untuk materi
selanjutnya).
Dengan melewati sebagian bagian, mungkin kita dapat membahas topik-topik lain
yang lebih lanjut.

\medskip

Ketergantungan umum bab-bab non-lampiran ditunjukkan dalam diagram berikut.
Biasanya saya menjalankan kuliah satu semester dengan menempuh bab
\ref{ch:complexplane}--\ref{ch:counting}, melewati versi homotopi Cauchy,
untuk menyelesaikan teori dasar fungsi holomorf, kemudian menuju
\ref{ch:montelriemann} (pemetaan Montel dan Riemann), serta beberapa bagian
dari \ref{ch:harmonic} (fungsi harmonik).
Ada beberapa topik tambahan untuk rencana yang berbeda, seperti
\ref{ch:weier} (faktorisasi Weierstrass), \ref{ch:runge} (Runge), dan
\ref{ch:analcont} (kelanjutan analitik).
\begin{equation*}
\begin{tikzcd}[cramped, row sep=small]
{\text{\Chdotref{ch:complexplane}}} \arrow[r] &
{\text{\Chdotref{ch:holanal}}} \arrow[r] &
{\text{\Chdotref{ch:cauchysimple}}} \arrow[r] &
{\text{\Chdotref{ch:log}}} \arrow[dl] \\
& & {\text{\Chdotref{ch:counting}}} \arrow[dl] \arrow[d]
\arrow[dr] \\
& {\text{\Chdotref{ch:montelriemann}}} \arrow[dl] \arrow[d] &
{\text{\Chdotref{ch:harmonic}}} &
{\text{\Chdotref{ch:weier}}}
\\
{\text{\Chdotref{ch:runge}}} &
{\text{\Chdotref{ch:analcont}}}
\end{tikzcd}
\end{equation*}
Satu-satunya alasan mengapa \ref{ch:runge} (Runge) bergantung pada
\ref{ch:montelriemann} (pemetaan Montel dan Riemann) adalah bahwa kita
membuktikan \lemmaref{lemma:patharoundK} (di sekitar setiap himpunan kompak
terdapat suatu siklus yang homolog dengan nol) sebagai contoh penerapan
pemetaan Riemann.
